\documentclass[11pt,a4paper]{article}

\usepackage[margin=25mm]{geometry}
\usepackage{amsmath,amssymb,amsthm}
\usepackage{fontspec}
\usepackage{xeCJK}
\usepackage{microtype}
\usepackage[hidelinks]{hyperref}

\setmainfont{Latin Modern Roman}
\setCJKmainfont[AutoFakeSlant=0.2]{Noto Serif CJK SC}
\setCJKmonofont{Noto Sans Mono CJK SC}
\setlength{\parindent}{0pt}
\setlength{\parskip}{0.65em}
\setlength{\jot}{7pt}

\newtheorem{theorem}{Theorem}
\newtheorem{corollary}{Corollary}

\title{\textbf{Fixed Points of a Random Permutation}\\
  \large Expected Value and Distribution}
\author{}
\date{}

\hypersetup{
  pdftitle={Fixed Points of a Random Permutation -- Revised English},
  pdfsubject={A self-contained report on derangements and fixed points}
}

\begin{document}

\newcommand{\EditionName}{Revised English edition}
\newcommand{\EditionSummary}{%
  This edition is a corrected, self-contained report based on the
  mathematical question in the source video.}
\hypersetup{pageanchor=false}
\begin{titlepage}
\thispagestyle{empty}
\centering

\vspace*{1.2cm}

{\Huge\bfseries Fixed Points of a Random Permutation\par}
\vspace{0.45cm}
{\Large Reproduction and Revised Editions\par}

\vspace{1.1cm}

\begin{minipage}{0.86\textwidth}
\raggedright

\textbf{Source.}
This document is an independent LaTeX reproduction of, and editorial
adaptation based on, the mathematical paper shown in the Bilibili video
\emph{【日常研究】摆烂小能手掀起的数学问题} (video ID
\texttt{BV1Fh411N71S}):
\begin{center}
  \href{https://www.bilibili.com/video/BV1Fh411N71S}
       {\nolinkurl{https://www.bilibili.com/video/BV1Fh411N71S}}
\end{center}

\textbf{Repository.}
The \LaTeX\ sources, compiled PDFs, and edition notes are published at:
\begin{center}
  \href{https://github.com/Ritsuka314/personal-website}
       {\nolinkurl{https://github.com/Ritsuka314/personal-website}}
\end{center}

\textbf{PDF editions.}
\begin{center}
  \small
  \href{https://opuscula.ritsuka.moe/euler-derangement-cloze/paper-en.pdf}
       {\texttt{paper-en.pdf}}
  \quad
  \href{https://opuscula.ritsuka.moe/euler-derangement-cloze/paper-en-revised.pdf}
       {\texttt{paper-en-revised.pdf}}
  \quad
  \href{https://opuscula.ritsuka.moe/euler-derangement-cloze/paper-zh.pdf}
       {\texttt{paper-zh.pdf}}
\end{center}

\textbf{Purpose.}
The reconstruction was prepared solely for entertainment and educational
purposes. It is unofficial, is not presented as an author-approved
transcript, and makes no claim of affiliation with the video creator or
platform.

\medskip
\textbf{The three editions.}
\begin{enumerate}
  \item \textbf{Reproduced English edition.}
    A close typesetting of the four-page English note visible in the video.
    It intentionally retains the source's nonstandard wording, omissions,
    notation, and induction argument.

  \item \textbf{Revised English edition.}
    A corrected, self-contained mathematical report. It defines the model
    precisely, proves the derangement formula by inclusion--exclusion, gives
    a unified probability distribution, and proves the expectation with
    indicator variables.

  \item \textbf{Revised Chinese edition.}
    A polished Chinese counterpart to the revised English report. It follows
    the same definitions, results, and proofs rather than translating the
    reproduced wording literally.
\end{enumerate}

\medskip
\noindent
\fbox{%
  \begin{minipage}{0.93\linewidth}
    \textbf{Edition in this file: \EditionName}\par
    \EditionSummary
  \end{minipage}%
}

\end{minipage}

\vfill
{\small Typeset with \LaTeX\ from the source video cited above.\par}

\end{titlepage}
\hypersetup{pageanchor=true}


\maketitle
\vspace{-2em}

\begin{abstract}
Let a permutation of \(n\) distinct objects be chosen uniformly at random,
and let \(X\) denote its number of fixed points. This report first proves
directly that \(\mathbb{E}[X]=1\) for every positive integer \(n\). It then
develops the derangement formula step by step using inclusion--exclusion and
derives the exact distribution of \(X\). The result models, for example, the
number of correct answers obtained when \(n\) distinct choices are assigned
uniformly at random to \(n\) distinct blanks.
\end{abstract}

\section{Problem setup}

For a positive integer \(n\), write
\[
  [n]=\{1,2,\ldots,n\}.
\]
Let \(\sigma:[n]\to[n]\) be a permutation chosen uniformly from the \(n!\)
elements of the symmetric group \(S_n\). An index \(r\in[n]\) is a
\emph{fixed point} of \(\sigma\) if \(\sigma(r)=r\). Define
\[
  X=\bigl|\{r\in[n]:\sigma(r)=r\}\bigr|.
\]
Our main question is the expected value of \(X\); afterward we determine its
entire probability mass function as a stronger result.

\section{Expected number of fixed points}

\begin{theorem}
For every \(n\geq 1\),
\[
  \mathbb{E}[X]=1.
\]
\end{theorem}

\begin{proof}
For each \(r\in[n]\), define the indicator variable
\[
  I_r=
  \begin{cases}
    1, & \sigma(r)=r,\\
    0, & \sigma(r)\ne r.
  \end{cases}
\]
Then \(X=\sum_{r=1}^{n}I_r\). Of the \(n!\) permutations, exactly
\((n-1)!\) fix a specified index \(r\), so
\[
  \Pr(\sigma(r)=r)=\frac{(n-1)!}{n!}=\frac{1}{n}
  \quad\text{and}\quad
  \mathbb{E}[I_r]=\frac{1}{n}.
\]
By linearity of expectation,
\[
  \mathbb{E}[X]
    =\sum_{r=1}^{n}\mathbb{E}[I_r]
    =n\cdot\frac{1}{n}
    =1.
\]
The indicator variables need not be independent for this argument.
\end{proof}

\textbf{Equivalent double-counting view.}
Count the pairs \((\sigma,r)\in S_n\times[n]\) for which
\(\sigma(r)=r\). For each of the \(n\) choices of \(r\), exactly
\((n-1)!\) permutations fix \(r\), so there are
\(n(n-1)!=n!\) such pairs. Dividing this total number of fixed-point
incidences by the \(n!\) permutations again shows that the average number of
fixed points is \(1\).

\section{Derangements and the exact distribution}

The expectation proof does not require the distribution of \(X\). To obtain
the stronger result \(\Pr(X=k)\) for every \(k\), we first count
permutations having no fixed points.

\subsection{Counting derangements}

A \emph{derangement} is a permutation with no fixed points. Let \(D_m\)
denote the number of derangements of \(m\) objects, with the convention
\(D_0=1\).

\begin{theorem}[Derangement formula]
For every integer \(m\geq 0\),
\[
  D_m=m!\sum_{j=0}^{m}\frac{(-1)^j}{j!}.
\]
Consequently, a uniformly random permutation of \(m\) objects is a
derangement with probability
\[
  \frac{D_m}{m!}=\sum_{j=0}^{m}\frac{(-1)^j}{j!}.
\]
\end{theorem}

\begin{proof}
For \(m=0\), the single empty permutation is a derangement, and the formula
gives \(D_0=1\). Now suppose \(m\geq 1\), and take the full set \(S_m\) of
all \(m!\) permutations as the universe. For each \(r\in[m]\), let
\[
  A_r=\{\sigma\in S_m:\sigma(r)=r\};
\]
thus \(A_r\) is the set of permutations that fix \(r\). A permutation is a
derangement precisely when it lies in none of \(A_1,\ldots,A_m\), so
\[
  D_m=\left|\bigcap_{r=1}^{m}\overline{A_r}\right|.
\]

Inclusion--exclusion begins with all \(m!\) permutations, subtracts those
that fix each specified element, adds back those subtracted twice because
they fix two specified elements, and continues with alternating signs:
\begin{align*}
  D_m
    ={}&|S_m|-\sum_{r=1}^{m}|A_r|
       +\sum_{1\leq r<s\leq m}|A_r\cap A_s|-\cdots\\
       &\quad+(-1)^m|A_1\cap\cdots\cap A_m|.
\end{align*}
To evaluate these terms, choose a subset \(J\subseteq[m]\) with
\(|J|=j\). If every element of \(J\) must be fixed, the remaining \(m-j\)
elements may be permuted arbitrarily. Hence
\[
  \left|\bigcap_{r\in J}A_r\right|=(m-j)!.
\]
When \(J=\varnothing\), the intersection is understood to be the whole
universe \(S_m\), so this statement also includes the \(j=0\) term.
There are \(\binom{m}{j}\) choices for \(J\). Grouping the
inclusion--exclusion terms according to \(j\) therefore gives
\begin{align*}
  D_m
    &=\sum_{j=0}^{m}(-1)^j\binom{m}{j}(m-j)!\\
    &=m!\sum_{j=0}^{m}\frac{(-1)^j}{j!}.
\end{align*}

For completeness, the alternating sum keeps exactly the desired
permutations. A permutation with \(q\geq1\) fixed points occurs in
\(\binom{q}{j}\) of the terms involving \(j\) prescribed fixed points, so
its net coefficient is
\[
  \sum_{j=0}^{q}(-1)^j\binom{q}{j}=(1-1)^q=0.
\]
A derangement has \(q=0\) and remains counted once. This also explains
directly why inclusion--exclusion works here.

Dividing by the total number \(m!\) of permutations proves the probability
formula.
\end{proof}

For example,
\[
  D_3
    =3!-\binom{3}{1}2!+\binom{3}{2}1!-\binom{3}{3}0!
    =6-6+3-1
    =2.
\]
In one-line notation, the two derangements are \([2,3,1]\) and
\([3,1,2]\).

When \(m\geq 2\), the terms for \(j=0\) and \(j=1\) cancel. Thus the same
probability may also be written as
\[
  \sum_{j=2}^{m}\frac{(-1)^j}{j!},
\]
which is the form used in the reproduced note.

\subsection{Distribution of the number of fixed points}

\begin{theorem}
For \(k=0,1,\ldots,n\),
\[
  \Pr(X=k)
    =\frac{\binom{n}{k}D_{n-k}}{n!}
    =\frac{1}{k!}\sum_{j=0}^{n-k}\frac{(-1)^j}{j!}.
\]
\end{theorem}

\begin{proof}
First choose the \(k\) indices that will be fixed; this can be done in
\(\binom{n}{k}\) ways. The remaining \(n-k\) indices must then be deranged,
which can be done in \(D_{n-k}\) ways. Hence exactly
\(\binom{n}{k}D_{n-k}\) permutations have \(k\) fixed points. Dividing by
\(n!\) and substituting the derangement formula yields
\begin{align*}
  \Pr(X=k)
    &=\frac{\binom{n}{k}D_{n-k}}{n!}\\
    &=\frac{1}{k!}
      \sum_{j=0}^{n-k}\frac{(-1)^j}{j!}.
\end{align*}
\end{proof}

The unified formula includes the two boundary cases:
\[
  \Pr(X=n)=\frac{1}{n!},
  \qquad
  \Pr(X=n-1)=0.
\]
The second identity holds because a permutation cannot move exactly one
element while fixing every other element.

\section{Conclusion}

The direct indicator argument shows that a uniformly random permutation has
one fixed point on average. The stronger counting argument gives the exact
distribution
\[
  \Pr(X=k)=\frac{\binom{n}{k}D_{n-k}}{n!},
\]
where the derangement number \(D_{n-k}\) is obtained by
inclusion--exclusion. Therefore, when \(n\) distinct answers are placed
uniformly at random into \(n\) distinct blanks, one answer is correct on
average. This is an expectation statement: an individual attempt may have
no correct answers, one correct answer, or several.

\end{document}
